% ============================================================
%  Глава 0 — Введение: кому, зачем, как читать
% ============================================================
\chapter{Зачем эта книжка}
\label{ch:intro}

\section{Кому это}

Этот документ написан для студента, который услышал «модель пространства
состояний» и «MPC» в курсе теории автоматического управления (ТАУ),
открыл код проекта \texttt{SAMURAI} --- и не понял, как одно стыкуется
с другим. Учебник Козлова даёт строгие формулы, но не показывает, как
эти формулы превращаются в команды моторам реального робота. Код проекта
работает и моторы крутятся, но «почему именно эти числа в матрице $\mat{A}$»
не объясняет.

\textbf{Книжка закрывает этот зазор.} Она ведёт читателя по реальной траектории
сценария \texttt{«проехать D метров вперёд»}: от нажатия кнопки в UI до
ШИМ-сигнала в моторе. На каждом шаге появляется ровно та математика,
которая нужна именно здесь, и тут же показывается, в каком файле проекта
она живёт.

\section{Что нужно знать заранее}

\begin{itemize}
  \item \textbf{Матан.} Производные, частные производные, ряд Тейлора первого
        порядка, неопределённый интеграл. Дифференциальные уравнения первого
        порядка с постоянными коэффициентами (типа $\dot{x} = -ax + bu$) ---
        не пугают, но если пугают, то мы их разберём заново.
  \item \textbf{Линейная алгебра.} Что такое матрица, умножение матриц,
        транспонирование, собственные значения. Знать слово «положительно
        определённая» полезно, но не обязательно --- введём.
  \item \textbf{Python на уровне чтения.} Уметь читать сигнатуру функции
        и срез массива \texttt{numpy}. Писать на нём не придётся ---
        главы про реализацию показывают код проекта и комментируют его.
\end{itemize}

Если матан и ЛА давно не открывал --- ничего страшного, нужные вещи
вспомним по ходу.

\section{Чего книжка \emph{не} объясняет}

Чтобы не разрастаться в энциклопедию, документ \textbf{намеренно не покрывает}
следующее:

\begin{itemize}
  \item Электронику и моторы: ШИМ, противо-ЭДС, H-мост. Считаем, что
        мотор --- это «чёрный ящик», который получает команду и крутится
        с некоторой инерцией. Подробнее --- в \filepath{latex\_doc/chapters/14\_pwm\_motor.tex}.
  \item Фильтрацию датчиков: EKF на IMU, слияние акселерометра и одометрии.
        Положение робота уже считаем «измеренным», источник --- глава
        \ref{ch:state}. Подробнее --- в \filepath{latex\_doc/chapters/01\_velocity\_ekf.tex}
        и далее.
  \item SLAM, навигацию A*, Pure Pursuit, патрулирование. Сценарий «D метров
        вперёд» не требует карты и глобального плана --- всё происходит
        в локальной системе координат робота.
\end{itemize}

\section{Как читать}

\textbf{Главы пронумерованы по сценарию использования, а не по сложности
математики.} В каждой главе сначала возникает \emph{вопрос} (например, «откуда
робот знает, на сколько он повернулся?»), потом \emph{интуиция} (на пальцах),
потом \emph{формула}, и в конце --- \emph{ссылка на файл и строку кода},
где это работает.

Если читаешь подряд, к главе~\ref{ch:mpc} ты уже будешь знать, что такое
$\mat{A}$, $\mat{B}$, $\mat{Q}$, $\mat{R}$, и MPC будет уже не «магия из
учебника», а «вот зачем мы их вводили».

Если читаешь нелинейно (только нужное), вот ориентир по типам коробок:

\begin{itemize}
  \item \textbf{Идея} --- интуиция, можно читать без формул.
  \item \textbf{Формула} --- ключевое уравнение, на которое потом ссылаемся.
  \item \textbf{Пример} --- численный пример, можно «прожать» в голове.
  \item \textbf{В коде} --- ссылка на файл/функцию в проекте.
  \item \textbf{Примечание} --- объяснение «в сторону», не обязательное.
  \item \textbf{Что если?} --- учебный «а что будет, если...». Полезно для
        интуиции.
  \item \textbf{Внимание} --- типичная ошибка, реальный подводный камень
        проекта.
\end{itemize}

\section{Что значит «МПС»}

\begin{ideabox}
В этом проекте \textbf{МПС = Модель Пространства Состояний}, по-английски
\emph{state-space model}. Это способ описать поведение робота уравнением
вида
\[
  \vect{x}_{k+1} = \mat{A}_d\,\vect{x}_k + \mat{B}_d\,\vect{u}_k,
\]
где $\vect{x}$ --- вектор состояния (что робот «есть в моменте»: где он,
куда смотрит, как быстро едет), а $\vect{u}$ --- вектор управления (команды
скорости, которые мы ему даём).

Аббревиатура «МПС» \textbf{не} означает «методы принятия решений» (про
Байеса, Гермейера, минимакс) --- если ты искал это, тебе нужна другая
ветка проекта.
\end{ideabox}

\section{Куда смотреть, если книжки мало}

\begin{itemize}
  \item \filepath{latex\_doc/control\_theory/} --- более формальный, сухой
        вариант того же материала. Хорошо для зачёта.
  \item \filepath{latex\_doc/main.tex} --- общий справочник по всей
        математике робота (фильтры, SLAM, кватернионы, A*) --- 15 глав.
  \item \filepath{docs/mps/} --- markdown-доки: \texttt{architecture.md},
        \texttt{api.md}, \texttt{scenario\_forward.md}. Описывают, как
        модуль устроен «снаружи».
  \item \filepath{docs/superpowers/specs/2026-05-05-mps-state-space-design.md}
        --- источник правды по проектным решениям feat/mps.
  \item Учебник: Козлов В.\,Н., Куприянов В.\,Е., Шашихин В.\,Н.
        «Теория автоматического управления», СПбГПУ, 2008.
\end{itemize}

\section{Обозначения}

\begin{center}
\begin{tabular}{ll}
\toprule
$\vect{x}(t) \in \Real^n$    & вектор состояния (непрерывное время)\\
$\vect{x}_k \in \Real^n$     & вектор состояния (дискретное время, шаг $k$)\\
$\vect{u}(t), \vect{u}_k$    & вектор управления, $r$ компонент\\
$\vect{y}(t), \vect{y}_k$    & вектор выхода (что «видит» наблюдатель), $m$ компонент\\
$\mat{A}, \mat{B}$           & матрицы непрерывной модели $\dot{\vect{x}} = \mat{A}\vect{x} + \mat{B}\vect{u}$\\
$\mat{A}_d, \mat{B}_d$       & матрицы дискретной модели $\vect{x}_{k+1} = \mat{A}_d\vect{x}_k + \mat{B}_d\vect{u}_k$\\
$\mat{C}, \mat{D}$           & матрицы выхода: $\vect{y} = \mat{C}\vect{x} + \mat{D}\vect{u}$\\
$\mat{Q}, \mat{R}$           & весовые матрицы критерия качества\\
$\mat{P}_f$                  & терминальный штраф (решение DARE)\\
$\mat{K}$                    & матрица обратной связи регулятора\\
$T_s$                        & период дискретизации (в проекте $0.02$~с)\\
$N$                          & горизонт прогноза MPC\\
$\vect{x}\T$                 & транспонирование\\
$\|\vect{x}\|_{\mat{Q}}^2$   & сокращение для $\vect{x}\T \mat{Q}\, \vect{x}$ (квадратичная форма)\\
$\bar{\vect{x}}, \bar{\vect{u}}$ & опорная точка линеаризации\\
$\Delta\vect{x}$             & отклонение от опорной точки: $\vect{x} - \bar{\vect{x}}$\\
\bottomrule
\end{tabular}
\end{center}

Переменные с тильдой (\texttt{\~{}x}) встречаются в коде --- это та же
$\vect{x}$, просто иначе оформлено в идентификаторах.

\medskip
Готов? Тогда поехали ---  смотрим, что вообще делает робот.
